\scanpage{212}
\mspage{211}
\hfill Tribu n° 59

SYSTÈMES de RACINES

\uline{Introduction}, (Rédaction Borel).

Ce chapitre est principalement consacré à l'étude de certains ensembles
finis de vecteurs d'un espace euclidien, appelés systèmes de racines, et des
groupes de déplacements d'un espace euclidien vectoriel ou affine engendré par
des réflexions. Son but est de fournir le cadre naturel dans lequel s'expriment
et se démontrent un grand nombre de propriétés importantes des groupes et algèbres
de Lie semi-simples sur $\underline{\underline{R}}$ ou sur $\underline{\underline{C}}$ et des groupes algébriques semi-simples
sur un corps quelconque. Bien que la rédaction ait été dictée par les besoins
des chapitres suivants de ce Livre, nous avons été amenés en plusieurs endroits
à élargir de diverses façons les hypothèses faites sur les groupes envisagés,
de manière à placer certains résultats dans leur domaine naturel de validité,
qui déborde parfois celui des applications que nous avons en vue. Le lecteur
uniquement intéressé à ces dernières peut concentrer son attention sur les
systèmes de racines, (§ 6), et sur leurs groupes de Weyl vectoriels ou affines.
Il peut se limiter, dans le § 1 aux réflexions d'un espace euclidien ou d'un
$\underline{\underline{Z}}$-module libre, dans le Théorème 1 du § 7 à l'implication (a) $\Longrightarrow$ (b) lorsque
$k=\underline{\underline{R}}$ et que $G$ est engendré par des réflexions, et envisager la discussion des
groupes de Coxeter faite au § 4 avant tout comme un auxiliaire permettant de
donner une démonstration unifiée des propriétés fondamentales des groupes de
déplacements linéaires ou affines engendrés par des réflexions $^{*}$.

\figure{\begin{tikzpicture}\draw[dashed] (0,0) -- (14,0);\end{tikzpicture}}

$(^{*})$ Pour rédiger ce paragraphe, nous avons largement utilisé des Notes de
Séminaire, (non publiées), de Jacques Tits.
